\documentclass[../main.tex]{subfiles}
\ifSubfilesClassLoaded{\setcounter{chapter}{5}}{}
\begin{document}

\chapter{Fungsi dan Modularitas}

\begin{subcpmk}
  \item Sub-CPMK 3.1: Mendefinisikan fungsi dengan parameter, nilai balik, dan ruang lingkup variabel lokal
  \item Sub-CPMK 3.2: Menggunakan argumen variadik (\texttt{*args}, \texttt{**kwargs}), ekspresi \textit{lambda}, dan anotasi tipe dasar pada fungsi
\end{subcpmk}

\SesiRPS{6--7}{3.1 dan 3.2}{$\sim$340 menit (dua sesi praktikum)}{CPMK-3}

\section{Sarana dan prasarana}

Python 3.10+, editor, terminal.

\section{Prosedur kerja}

\begin{enumerate}[leftmargin=*]
  \item Definisikan fungsi dengan parameter dan nilai balik; uji dari REPL.
  \item Tambahkan \texttt{*args}/\texttt{**kwargs}, \textit{lambda}, dan anotasi tipe sesuai contoh modul.
  \item Refaktor program monolitik menjadi fungsi-fungsi kecil.
\end{enumerate}

\section{Definisi Fungsi}

Fungsi memecah program menjadi bagian yang dapat digunakan kembali. Didefinisikan dengan \keyword{def}:

\begin{lstlisting}[caption=Fungsi sederhana (Tutorial Python ss4.8)]
def kuadrat(x):
    """Mengembalikan kuadrat dari x."""
    return x * x

def salam(nama):
    """Mencetak salam kepada nama."""
    print("Halo,", nama)

print(kuadrat(5))   # 25
salam("Aulia")      # Halo, Aulia
\end{lstlisting}

Pernyataan \textbf{pertama} dari badan fungsi bisa berupa \textit{docstring} (dokumentasi fungsi dalam string literal). Gunakan \texttt{help(kuadrat)} di REPL untuk membacanya.

\section{Parameter dan Nilai Balik}

Fungsi dapat mengembalikan beberapa nilai sekaligus dalam bentuk tuple:

\begin{lstlisting}[caption=Nilai balik berganda]
def statistik(data):
    """Mengembalikan (minimum, maksimum, rata-rata)."""
    return min(data), max(data), sum(data) / len(data)

nilai = [75, 90, 60, 85, 78]
mn, mx, rata = statistik(nilai)
print(f"Min={mn}, Max={mx}, Rata={rata:.2f}")
\end{lstlisting}

\section{Nilai Argumen Default}

Parameter dapat memiliki nilai \textit{default} yang digunakan bila argumen tidak diberikan saat pemanggilan (Tutorial Python §4.9.1):

\begin{lstlisting}[caption=Nilai argumen default]
def hitung_pangkat(basis, eksponen=2):
    """Menghitung basis ** eksponen (default: kuadrat)."""
    return basis ** eksponen

print(hitung_pangkat(3))      # 9  (eksponen=2 default)
print(hitung_pangkat(3, 3))   # 27
print(hitung_pangkat(2, 10))  # 1024
\end{lstlisting}

\begin{catatan}
Nilai default dievaluasi \textbf{sekali} saat definisi fungsi, bukan setiap kali dipanggil. Jangan gunakan objek \textit{mutable} (seperti \texttt{list}) sebagai nilai default karena bisa menyebabkan bug halus.

\textbf{Salah:} \texttt{def tambah(data, daftar=[]):} --- daftar yang sama dipakai ulang!\\
\textbf{Benar:} \texttt{def tambah(data, daftar=None): if daftar is None: daftar = []}
\end{catatan}

\section{Keyword Arguments}

Saat memanggil fungsi, argumen dapat dilewatkan dengan nama parameter secara eksplisit (Tutorial Python §4.9.2):

\begin{lstlisting}[caption=Keyword arguments]
def cetak_info(nama, umur, kota="Jakarta"):
    print(f"{nama}, {umur} tahun, dari {kota}")

# Pemanggilan dengan keyword args
cetak_info("Budi", 20)
cetak_info("Sari", umur=22, kota="Bandung")
cetak_info(kota="Surabaya", nama="Andi", umur=19)
\end{lstlisting}

\section{Argumen Tak Tentu: \texttt{*args} dan \texttt{**kwargs}}

Untuk fungsi yang menerima jumlah argumen tak tentu (Tutorial Python §4.9.4):

\begin{lstlisting}[caption=*args dan **kwargs]
def jumlahkan(*angka):
    """Menjumlahkan semua argumen yang diberikan."""
    total = 0
    for n in angka:
        total += n
    return total

print(jumlahkan(1, 2, 3))           # 6
print(jumlahkan(10, 20, 30, 40))    # 100

def cetak_profil(**info):
    """Mencetak pasangan kunci-nilai."""
    for kunci, nilai in info.items():
        print(f"  {kunci}: {nilai}")

cetak_profil(nama="Rini", nim="12345", ipk=3.8)
\end{lstlisting}

\section{Ekspresi Lambda}

Lambda adalah fungsi kecil anonim (tanpa nama) yang ditulis dalam satu baris (Tutorial Python §4.9.6):

\begin{lstlisting}[caption=Lambda expressions]
# Fungsi biasa
def kali_dua(x):
    return x * 2

# Ekuivalen lambda
kali_dua = lambda x: x * 2
print(kali_dua(5))   # 10

# Lambda berguna sebagai argumen ke fungsi lain
nilai = [("Andi", 85), ("Budi", 72), ("Citra", 90)]
nilai.sort(key=lambda item: item[1])   # sort berdasarkan nilai
print(nilai)

# sorted dengan lambda
angka = [5, 1, 8, 3, 9, 2]
terurut = sorted(angka, key=lambda x: -x)  # urutan menurun
print(terurut)   # [9, 8, 5, 3, 2, 1]
\end{lstlisting}

\begin{konsep}
Lambda terbatas pada \textbf{satu ekspresi}. Untuk logika yang lebih kompleks, gunakan fungsi biasa dengan \keyword{def}. Lambda paling berguna saat dipakai sebagai argumen pada fungsi seperti \texttt{sorted()}, \texttt{map()}, \texttt{filter()}.
\end{konsep}

\section{Anotasi Fungsi (Type Hints)}

Python mendukung anotasi tipe untuk mendokumentasikan tipe parameter dan nilai balik (Tutorial Python §4.9.8). Anotasi tidak mempengaruhi eksekusi program, tetapi membantu pembaca dan alat analisis statis:

\begin{lstlisting}[caption=Anotasi fungsi / type hints]
def hitung_luas(panjang: float, lebar: float) -> float:
    """Menghitung luas persegi panjang."""
    return panjang * lebar

def sapa(nama: str, formal: bool = False) -> str:
    """Membuat kalimat salam."""
    if formal:
        return f"Selamat pagi, {nama}."
    return f"Halo, {nama}!"

print(hitung_luas(5.0, 3.0))     # 15.0
print(sapa("Budi", formal=True))  # Selamat pagi, Budi.
\end{lstlisting}

\section{Ruang Lingkup Variabel}

Variabel yang diberi nilai di dalam fungsi bersifat \textit{lokal} kecuali dinyatakan lain dengan \keyword{global} atau \keyword{nonlocal}. Penggunaan \keyword{global} jarang disarankan karena membuat kode sulit dilacak:

\begin{lstlisting}[caption=Ruang lingkup lokal]
x = 10   # variabel global

def ubah_lokal():
    x = 99   # variabel lokal, berbeda dari global x
    print("Dalam fungsi:", x)

ubah_lokal()
print("Di luar fungsi:", x)   # masih 10
\end{lstlisting}

\section{Studi Kasus dan Contoh Kode}

Berikut adalah contoh-contoh program Python yang mengimplementasikan konsep-konsep pada modul ini.

\subsection{Kode: \texttt{faktorial.py}}
\begin{lstlisting}[language=Python,caption={\texttt{faktorial.py}}]
"""Faktorial rekursif (rumus Pascal sumber salah; di Python dipakai n * fakt(n-1))."""


def fakt(n: int) -> int:
    if n == 0:
        return 1
    return n * fakt(n - 1)


def main() -> None:
    n = int(input("Masukkan n : "))
    print(n, "! = ", fakt(n), sep="")


if __name__ == "__main__":
    main()
\end{lstlisting}

\subsection{Kode: \texttt{prosedur\_upah\_latihan.py}}
\begin{lstlisting}[language=Python,caption={\texttt{prosedur\_upah\_latihan.py}}]
"""Latihan prosedur/fungsi upah - modul Pascal tidak lengkap; stub di Python."""


def input_pegawai(nama: str, gol: str, jml_anak: int) -> None:
    # TODO: definisikan sendiri (parameter seharusnya diisi via input)
    raise NotImplementedError("Lengkapi input_pegawai.")


def upah_kotor(gol: str) -> float:
    # TODO: definisikan sendiri
    raise NotImplementedError("Lengkapi upah_kotor.")


def persen_tunjangan(jml_anak: int) -> float:
    # TODO: definisikan sendiri
    raise NotImplementedError("Lengkapi persen_tunjangan.")


def main() -> None:
    print(
        "Modul latihan: lengkapi deklarasi variabel dan fungsi seperti di modul Pascal."
    )
    try:
        # Contoh alur yang dimaksud Pascal (setelah variabel didefinisikan):
        # input(...); upah_bersih = upah_kotor(gol) - (upah_kotor(gol) * persen_tunjangan(jml_anak))
        raise NotImplementedError("Definisikan Nama, Gol, jml_anak lalu panggil fungsi.")
    except NotImplementedError as e:
        print(e)


if __name__ == "__main__":
    main()
\end{lstlisting}

\subsection{Kode: \texttt{tabel\_latihan.py}}
\begin{lstlisting}[language=Python,caption={\texttt{tabel\_latihan.py}}]
"""Latihan tabel pangkat - lengkapi fpangkat dan cetak (stub, setara tabel_latihan.pas)."""


def fpangkat(dipangkatkan: float, pangkat: int) -> float:
    # TODO: lengkapi (dipakai di dalam cetak bila perlu)
    raise NotImplementedError("Lengkapi fpangkat.")


def cetak(x: int) -> None:
    # TODO: lengkapi isi tabel seperti tabel_pangkat_x
    print("  x      1/x        x^2      x^3")
    print("-----------------------------------")
    raise NotImplementedError("Lengkapi badan cetak(x).")


def main() -> None:
    x = int(input("Masukkan banyaknya data : "))
    try:
        cetak(x)
    except NotImplementedError as e:
        print(e)
    input("Tekan Enter untuk melanjutkan...")


if __name__ == "__main__":
    main()
\end{lstlisting}

\section{Referensi sesi}

\begin{itemize}[leftmargin=*]
  \item \textbf{[12]} Sweigart, A. \textit{Automate the Boring Stuff with Python}.
  \item \textbf{[14]} Pilgrim, M. \textit{Dive Into Python 3}.
  \item \textbf{[8, 9]} \textit{The Python Tutorial}, definisi fungsi.
\end{itemize}

\begin{aktivitas}
  \item Tulis fungsi \texttt{diskon(harga, persen=10)} yang mengembalikan harga setelah diskon (dengan nilai default 10\%).
  \item Tulis fungsi \texttt{faktorial(n)} rekursif atau iteratif; bandingkan keterbacaan keduanya.
  \item Pecah program hitung statistik (min, max, rata-rata) dari daftar menjadi beberapa fungsi terpisah.
  \item Buat fungsi \texttt{gabung(*kata, pemisah=" ")} yang menggabungkan argumen kata dengan pemisah yang dapat dikustomisasi.
  \item Gunakan \texttt{sorted()} dengan lambda untuk mengurutkan list of dict berdasarkan field tertentu.
  \item Tambahkan type hints ke semua fungsi yang sudah Anda buat.
\end{aktivitas}

\begin{latihan}
  \item Apa perbedaan fungsi yang mengembalikan nilai dengan yang hanya mencetak?
  \item Kapan docstring berguna?
  \item Jelaskan kesalahan umum \textit{UnboundLocalError} pada variabel lokal/global.
  \item Apa keuntungan menggunakan keyword arguments? Berikan contoh kasus nyata.
  \item Kapan lambda cocok digunakan? Kapan sebaiknya menggunakan \keyword{def} biasa?
  \item Refleksi: bagaimana fungsi membantu pengujian program?
\end{latihan}

\begin{asesmen}
\textbf{Pilihan ganda}

\begin{enumerate}
  \item Pernyataan yang benar tentang \texttt{return}:
  \begin{enumerate}
    \item Wajib ada di setiap fungsi
    \item Menghentikan eksekusi fungsi dan dapat membawa nilai
    \item Hanya boleh satu per fungsi
    \item Tidak boleh digunakan di fungsi rekursif
  \end{enumerate}
  \item Definisi fungsi \texttt{def f(a, b=5):} berarti:
  \begin{enumerate}
    \item \texttt{b} selalu bernilai 5
    \item \texttt{b} bernilai 5 jika tidak diberikan saat pemanggilan
    \item \texttt{a} bersifat opsional
    \item Fungsi hanya dapat dipanggil dengan dua argumen
  \end{enumerate}
  \item Ekspresi lambda \texttt{lambda x, y: x + y} ekuivalen dengan:
  \begin{enumerate}
    \item \texttt{def f(x): return x + y}
    \item \texttt{def f(x, y): return x + y}
    \item \texttt{def f(): return x + y}
    \item \texttt{x + y}
  \end{enumerate}
\end{enumerate}

\textbf{Tugas}: tiga fungsi berikut dipanggil dari menu teks interaktif (tambahkan type hints dan docstring):\newline
\texttt{luas\_persegi(s: float) -> float},\newline
\texttt{luas\_segitiga(a: float, t: float) -> float},\newline
\texttt{luas\_lingkaran(r: float) -> float}.
\end{asesmen}

\begin{checklist}
  \item Saya dapat mendefinisikan fungsi dengan parameter dan \texttt{return}
  \item Saya dapat menggunakan nilai argumen default
  \item Saya dapat memanggil fungsi dengan keyword arguments
  \item Saya dapat mendefinisikan fungsi dengan \texttt{*args} dan \texttt{**kwargs}
  \item Saya dapat menulis ekspresi lambda sederhana
  \item Saya dapat menambahkan anotasi tipe (type hints)
  \item Saya dapat menjelaskan ruang lingkup variabel lokal
  \item Saya dapat menulis docstring yang informatif
\end{checklist}

\begin{rangkuman}
Fungsi meningkatkan keterbacaan, pengujian, dan penggunaan ulang kode. Python mendukung nilai default, keyword arguments, variadic args, lambda, dan type hints untuk membuat fungsi yang fleksibel dan expressif.

\textbf{Referensi:} {\small Python Tutorial §4.8--4.9 \url{https://docs.python.org/3/tutorial/controlflow.html}}

\textbf{Kata kunci:} \code{def}, \code{return}, \code{lambda}, \code{*args}, \code{**kwargs}, \code{docstring}, \code{type hints}, ruang lingkup lokal
\end{rangkuman}

\end{document}
